% Source: Weinberg GR, physical PDF pp. 347--352
%         (printed pp. 325--330).
% Coverage: Section 11.5, Supermassive Stars; equations
%           (11.5.1)--(11.5.18).
% Figures/tables/footnotes: no figures or tables; reference callouts 23--24.
% Status: source-reviewed and compile-clean.
% Uncertainties: none.

\subsection{Supermassive Stars}\label{sec:11.5}

We now turn to a different kind of
``star,''\hyperref[ch11-ref-23]{\textsuperscript{23}} in which general
relativity enters in quite a different way. Let us consider a Newtonian star
that is supported by the pressure of radiation rather than of matter; the
conditions under which this occurs will be discovered as we go along. Let us
also assume that the star is in convective equilibrium (see
Section~\ref{sec:11.1}) and has uniform chemical composition. Radiation has
an energy density \(e=3p\), so this star will be a polytrope with
\(\gamma=4/3\), that is,
\begin{equation}
  p=K\rho^{4/3}.
  \tag{11.5.1}\label{eq:11.5.1}
\end{equation}
The radiation pressure is given by the Stefan--Boltzmann law,
\begin{equation}
  p_r=\frac{\pi^2(kT)^4}{45\hbar^3},
  \tag{11.5.2}\label{eq:11.5.2}
\end{equation}
so with \(p\simeq p_r\), the temperature is given by
\begin{equation}
  kT
  =
  \left(\frac{45\hbar^3K}{\pi^2}\right)^{1/4}\rho^{1/3}.
  \tag{11.5.3}\label{eq:11.5.3}
\end{equation}
The pressure of matter here is given by the ideal-gas law:
\begin{equation}
  p_m=\rho\frac{kT}{\bar m},
  \tag{11.5.4}\label{eq:11.5.4}
\end{equation}
where \(\bar m\) is the mean mass of the gas particles. Thus the ratio of
matter to radiation pressure is
\begin{equation}
  \beta
  \equiv
  \frac{p_m}{p_r}
  =
  \frac{45\hbar^3}{\pi^2\bar m}\frac{\rho}{(kT)^3}
  =
  \frac1{\bar m}
  \left(\frac{45\hbar^3}{\pi^2K^3}\right)^{1/4}.
  \tag{11.5.5}\label{eq:11.5.5}
\end{equation}
This is a constant throughout the star, so we can use \(\beta\) instead of
\(K\) (or the entropy per nucleon, on which they both depend) to define the
equation of state, writing
\begin{equation}
  K
  =
  \left(\frac{45\hbar^3}{\pi^2\bar m^4\beta^4}\right)^{1/3}.
  \tag{11.5.6}\label{eq:11.5.6}
\end{equation}
The mass of a polytrope with \(\gamma=4/3\) is given by
Eq.~\eqref{eq:11.3.17} and Table~\ref{tab:11.1} as
\begin{equation}
  M
  =
  4\pi(2.01824)\left(\frac K{\pi G}\right)^{3/2},
  \tag{11.5.7}\label{eq:11.5.7}
\end{equation}
and, using Eq.~\eqref{eq:11.5.6}, this becomes
\begin{equation}
  \begin{split}
  M
  &=
  \frac{12\sqrt5}{\pi^{3/2}}(2.01824)
  \frac{\hbar^{3/2}}{\bar m^2G^{3/2}\beta^2}\\
  &=
  18M_\odot\left(\frac{m_N}{\bar m}\right)^2\beta^{-2}.
  \end{split}
  \tag{11.5.8}\label{eq:11.5.8}
\end{equation}
For ionized hydrogen at temperatures between \(10^{5.5}\ \mathrm{K}\) and
\(10^{10}\ \mathrm{K}\), \(\bar m\) is the average of the proton and the
electron mass, so \(\bar m\simeq m_N/2\). Thus in this case the condition for
radiation pressure to dominate material pressure by, say, a factor \(10\), is
that \(M\gtrsim7200M_\odot\). No such supermassive star has been definitely
observed, but they have been considered as possible sites for the production
of radiant energy through gravitational
collapse.\hyperref[ch11-ref-23]{\textsuperscript{23}}

The structure of a supermassive star is entirely determined by the equations
for a Newtonian polytrope with \(\gamma\simeq4/3\). In particular,
Eq.~\eqref{eq:11.3.16} gives the radius of the star as
\[
  R=6.89685\left(\frac K{\pi G}\right)^{1/2}\rho(0)^{-1/3},
\]
and, using Eq.~\eqref{eq:11.5.6}, this is
\begin{equation}
  R
  =
  \left(\frac{45}{\pi^5}\right)^{1/6}(6.89685)
  \frac{\hbar^{1/2}}
  {\bar m^{2/3}G^{1/2}\beta^{2/3}}
  \rho(0)^{-1/3}.
  \tag{11.5.9}\label{eq:11.5.9}
\end{equation}
The radius is restricted by our assumption that the rest-mass energy of the
star be much greater than its radiant energy, and, \emph{a fortiori}, than the
thermal energy of its matter. This condition reads
\[
  \frac{\pi^2(kT)^4}{15\hbar^3}\ll\rho,
\]
or, using Eqs.~\eqref{eq:11.5.3} and \eqref{eq:11.5.6},
\begin{equation}
  \rho\ll
  \frac{\pi^2}{1215}\frac{\beta^4\bar m^4}{\hbar^3}.
  \tag{11.5.10}\label{eq:11.5.10}
\end{equation}
The density \(\rho\) is greatest at the center, so this can be regarded as a
condition on \(\rho(0)\); using Eqs.~\eqref{eq:11.5.8} and
\eqref{eq:11.5.9} to express \(\beta\) and \(\rho(0)\) in terms of \(M\) and
\(R\), the condition \eqref{eq:11.5.10} becomes
\begin{equation}
  \frac{MG}{R}
  \ll
  \frac43\left(\frac{2.01824}{6.89685}\right)
  =
  0.39.
  \tag{11.5.11}\label{eq:11.5.11}
\end{equation}
This is essentially equivalent to the statement that the gravitational
potential is small, which also was assumed. For \(M=10^4M_\odot\),
Eq.~\eqref{eq:11.5.11} requires that \(R\gg4\times10^4\ \mathrm{km}\).

Although we do not need general relativity to understand the structure of
these supermassive stars, we shall need it to settle the question of stability.
A polytrope with \(\gamma=4/3\) is trembling between stability and
instability, so it is necessary to take into account the small effects of the
matter pressure and of general relativity, which play no appreciable role in
structure calculations.

We use Theorem~1 of Section~\ref{sec:11.2}, which tells us that the
transition from stability to instability will occur at a value of \(\rho(0)\)
for which the internal energy \(E\) is stationary. To calculate \(E\), we use
Eqs.~\eqref{eq:11.1.29}--\eqref{eq:11.1.31}, which to first order in
\(GM/R\) give
\begin{equation}
  \begin{aligned}
  E\simeq{}&
  \int_0^R4\pi r^2e(r)\,\dd r
  +\int_0^R4\pi Gr\mathcal{M}(r)e(r)\,\dd r\\
  &-\int_0^R4\pi Gr\mathcal{M}(r)\rho(r)\,\dd r
  -\int_0^R6\pi G^2\mathcal{M}^2(r)\rho(r)\,\dd r.
  \end{aligned}
  \tag{11.5.12}\label{eq:11.5.12}
\end{equation}
The internal energy density \(e\) is
\[
  e
  =
  \frac{\pi^2(kT)^4}{15\hbar^3}
  +\frac1{\Gamma-1}\frac{\rho kT}{\bar m}
  =
  3p_r\left[1+\frac{\beta}{3(\Gamma-1)}\right],
\]
where \(\Gamma\) is the specific-heat ratio of the matter. (For ionized
hydrogen, \(\Gamma=5/3\).) The total pressure is
\[
  p=p_r+p_m=p_r(1+\beta).
\]
Therefore, to first order in the small parameter \(\beta\), the ratio of
energy density to pressure is given by
\begin{equation}
  e
  \simeq
  3p
  \left[
    1-\frac{3\Gamma-4}{3(\Gamma-1)}\beta
    +\order(\beta^2)
  \right].
  \tag{11.5.13}\label{eq:11.5.13}
\end{equation}
The small correction of order \(\beta\) can be ignored in the second term in
Eq.~\eqref{eq:11.5.12}, which is already smaller than the first term by a
factor of order \(GM/R\); but it must be kept in the large first term, and
therefore
\begin{equation}
  \begin{aligned}
  E\simeq{}&
  \left[1-\frac{3\Gamma-4}{3(\Gamma-1)}\beta\right]
  \int_0^R12\pi r^2p(r)\,\dd r\\
  &+\int_0^R12\pi Gr\mathcal{M}(r)p(r)\,\dd r\\
  &-\int_0^R4\pi Gr\mathcal{M}(r)\rho(r)\,\dd r\\
  &-\int_0^R6\pi G^2\mathcal{M}^2(r)\rho(r)\,\dd r-\cdots.
  \end{aligned}
  \tag{11.5.14}\label{eq:11.5.14}
\end{equation}
The first integral can be rewritten by integrating by parts:
\[
  \int_0^R12\pi r^2p(r)\,\dd r
  =
  \int_0^Rp(r)\,\dd(4\pi r^3)
  =
  -\int_0^R4\pi r^3p'(r)\,\dd r.
\]
To calculate \(p'(r)\), we expand the fundamental equation
\eqref{eq:11.1.13} to first order in \(GM/R\):
\[
  -r^2p'(r)
  \simeq
  G\mathcal{M}(r)\rho(r)
  \left[
    1+\frac{p(r)}{\rho(r)}
    +\frac{4\pi r^3p(r)}{\mathcal{M}(r)}
    +\frac{2G\mathcal{M}(r)}r
  \right],
\]
so
\[
  \begin{aligned}
  \int_0^R12\pi r^2p(r)\,\dd r
  \simeq{}&
  \int_0^R4\pi Gr\mathcal{M}(r)\rho(r)\,\dd r\\
  &+\int_0^R4\pi Gr\mathcal{M}(r)p(r)\,\dd r\\
  &+\int_0^R16\pi^2Gr^4\rho(r)p(r)\,\dd r\\
  &+\int_0^R8\pi G^2\rho(r)\mathcal{M}^2(r)\,\dd r.
  \end{aligned}
\]
The \(\beta\)-correction needs to be kept in only the first term, which is
larger than the others by a factor of order \(R/(MG)\), so to first order in
\(\beta\) and \(GM/R\), Eq.~\eqref{eq:11.5.14} reads
\begin{equation}
  \begin{aligned}
  E\simeq{}&
  -\frac{3\Gamma-4}{3(\Gamma-1)}\beta
  \int_0^R4\pi Gr\mathcal{M}(r)\rho(r)\,\dd r\\
  &+\int_0^R16\pi Gr\mathcal{M}(r)p(r)\,\dd r\\
  &+\int_0^R16\pi^2Gr^4\rho(r)p(r)\,\dd r\\
  &+\int_0^R2\pi G^2\mathcal{M}^2(r)\rho(r)\,\dd r.
  \end{aligned}
  \tag{11.5.15}\label{eq:11.5.15}
\end{equation}
Now every term is small, so they can all be evaluated using for \(\rho\),
\(p\), and \(\mathcal{M}\) the values obtained by solving the Newtonian
equation
\[
  -r^2p'(r)\simeq G\mathcal{M}(r)\rho(r)
\]
for a Newtonian polytrope with \(\gamma=4/3\). In particular, the first
integral in Eq.~\eqref{eq:11.5.15} is given by setting \(\gamma=4/3\) in
Eq.~\eqref{eq:11.3.24},
\[
  \int_0^R4\pi Gr\mathcal{M}(r)\rho(r)\,\dd r
  =
  -V
  =
  \frac{3GM^2}{2R},
\]
whereas an integration by parts lets us write the third term as
\[
  \begin{aligned}
  \int_0^R16\pi^2Gr^4\rho(r)p(r)\,\dd r
  &=
  \int_0^R4\pi Gr^2p(r)\,\dd\mathcal{M}(r)\\
  &=
  -\int_0^R4\pi Gr^2p'(r)\mathcal{M}(r)\,\dd r
  -\int_0^R8\pi Grp(r)\mathcal{M}(r)\,\dd r\\
  &=
  \int_0^R4\pi G^2\mathcal{M}^2(r)\rho(r)\,\dd r
  -\int_0^R8\pi Grp(r)\mathcal{M}(r)\,\dd r.
  \end{aligned}
\]
Equation~\eqref{eq:11.5.15} now reads
\[
  \begin{aligned}
  E\simeq{}&
  -\frac{3\Gamma-4}{2(\Gamma-1)}\beta\frac{GM^2}{R}\\
  &+\int_0^R8\pi Gr\mathcal{M}(r)p(r)\,\dd r
  +\int_0^R6\pi G^2\mathcal{M}^2(r)\rho(r)\,\dd r.
  \end{aligned}
\]
The last two integrals may be calculated in terms of the Lane--Emden function
\(\theta(\xi)\) for \(\gamma=4/3\):
\[
  \begin{aligned}
  \int_0^R6\pi G^2\mathcal{M}^2(r)\rho(r)\,\dd r
  &=
  \frac{6K^{7/2}\rho(0)^{2/3}}{\pi^{5/2}G^{3/2}}
  \int_0^{\xi_1}\xi^4\theta'^2(\xi)\theta^3(\xi)\,\dd\xi,\\
  \int_0^R8\pi G\mathcal{M}(r)p(r)r\,\dd r
  &=
  \frac{8K^{7/2}\rho(0)^{2/3}}{\pi^{3/2}G^{3/2}}
  \int_0^{\xi_1}\xi^3|\theta'(\xi)|\theta^4(\xi)\,\dd\xi,
  \end{aligned}
\]
whereas \(K\) and \(\rho(0)\) can be expressed in terms of \(M\) and \(R\)
by using Eqs.~\eqref{eq:11.3.16} and \eqref{eq:11.3.17}, yielding
\[
  \frac{K^{7/2}\rho(0)^{2/3}}{G^{3/2}}
  =
  \frac{\sqrt\pi}{64\xi_1^4|\theta'(\xi_1)|^3}
  \frac{GM^2}{R}.
\]
A numerical integration
gives\hyperref[ch11-ref-24]{\textsuperscript{24}}
\[
  \frac1{8\pi\xi_1^4|\theta'(\xi_1)|^3}
  \left\{
    \int_0^{\xi_1}\xi^3|\theta'(\xi)|\theta^4(\xi)\,\dd\xi
    +\frac3{4\pi}
    \int_0^{\xi_1}\xi^4\theta'^2(\xi)\theta^3(\xi)\,\dd\xi
  \right\}
  =
  5.1,
\]
so, putting this all together, we have at last
\begin{equation}
  E
  \simeq
  -\frac{3\Gamma-4}{2(\Gamma-1)}
  \beta\frac{GM^2}{R}
  +5.1\frac{G^2M^3}{R^2}.
  \tag{11.5.16}\label{eq:11.5.16}
\end{equation}

The star is certainly stable when \(R\) is so large that general relativity
can be neglected, for then the star behaves like a Newtonian polytrope with
\[
  \gamma
  =
  1+\frac pe
  \simeq
  \frac43+\frac{3\Gamma-4}{9(\Gamma-1)}\beta
  >
  \frac43.
\]
[See Eq.~\eqref{eq:11.5.13}.] The transition from stability to instability
will occur when \(R\) decreases to a value where
\[
  \frac{\dd E}{\dd R}
  =
  \frac{\partial E}{\partial\rho(0)}
  \frac{\partial\rho(0)}{\partial R}
  =
  0.
\]
The derivative must be taken with constant entropy per nucleon, and hence in
this case with \(\beta\) fixed and \(M\) fixed. [See
Eqs.~\eqref{eq:11.5.6} and \eqref{eq:11.5.7}.] Thus the minimum radius for
stability is
\begin{equation}
  R_{\min}
  \simeq
  \frac{20.4(\Gamma-1)}{3\Gamma-4}\frac{GM}{\beta}.
  \tag{11.5.17}\label{eq:11.5.17}
\end{equation}
The maximum energy that can be released by letting the star shrink slowly
(through radiation at its surface) to this minimum stable radius is
\begin{equation}
  -E(R_{\min})
  \simeq
  \frac{(3\Gamma-4)^2\beta^2M}{81.6(\Gamma-1)^2}.
  \tag{11.5.18}\label{eq:11.5.18}
\end{equation}
For instance, a star with \(\beta=0.1\) will have
\(M\simeq7200M_\odot\); if \(\Gamma=5/3\), then the minimum radius is
\(1.45\times10^6\ \mathrm{km}\), and the fraction of its rest mass that can
be released by assembling the star is \(0.03\%\). The maximum value of the
surface potential \(MG/R\) for \(\Gamma=5/3\) is \(0.0735\beta\), well under
the limit \eqref{eq:11.5.11}.
